% Section 3.1: 数据来源与任务定义

本章为 CodeGeeX4-ALL-9B\cite{thudm2024codegeex4all9b} 的 LoRA 微调与 pass@$k$\cite{chen2021evaluating} 评测构建统一语料。原始来源选用公开的 KodCode\cite{xu2025kodcode}——其条目天然附带题目、参考解答与针对性单测，契合「生成—执行—判定」闭环。

\paragraph{（1）原始字段}
KodCode 以 Parquet 列式存储，每条记录含四个字段：
\begin{itemize}
    \item \texttt{question}：自然语言题目描述；
    \item \texttt{solution}：Python 参考解答；
    \item \texttt{test}：pytest 风格单测脚本；
    \item \texttt{test\_info}：被测函数签名（\texttt{function\_declaration}、\texttt{function\_name}、\texttt{parameter\_list}）。
\end{itemize}

\paragraph{（2）任务形式化与评测指标}
给定题目 $q$，模型需输出可直接执行的 Python 解答 $\hat{s}$；将 $\hat{s}$ 与官方 $t$ 组合在隔离子进程中运行，全部断言通过即记一次成功。令 $n$ 为采样次数、$c$ 为成功次数，
\[
\text{pass@}k \;=\; \mathbb{E}_{q}\!\left[\,1 - \binom{n-c}{k}\Big/\binom{n}{k}\,\right].
\]

\paragraph{（3）数据流水线总览}
从原始 Parquet 到训练/评测 JSONL 的五阶段流程见图~\ref{fig:ch3-pipeline}，最终产出训练集 $\approx$9000 条、评测集 1000 条。

\begin{figure}[!htb]
  \centering
  \begin{tikzpicture}[
    node distance=6mm and 5mm,
    stage/.style={draw, rounded corners=2pt, align=center, minimum height=9mm, minimum width=22mm, font=\small},
    tag/.style={font=\footnotesize\itshape, align=center},
    arr/.style={-{Stealth[length=2mm]}, thick},
  ]
    \node[stage] (s1) {原始\\ Parquet};
    \node[stage, right=of s1] (s2) {格式规范化\\ + 单测校验};
    \node[stage, right=of s2] (s3) {思维步骤\\ 标注};
    \node[stage, right=of s3] (s4) {算法标签\\ 标注};
    \node[stage, right=of s4] (s5) {顺序划分\\ 9{:}1};

    \draw[arr] (s1) -- (s2);
    \draw[arr] (s2) -- (s3);
    \draw[arr] (s3) -- (s4);
    \draw[arr] (s4) -- (s5);

    \node[tag, below=1mm of s1] {KodCode.parquet};
    \node[tag, below=1mm of s2] {KodCode.jsonl\\ ($\le$10000)};
    \node[tag, below=1mm of s3] {+\texttt{thought\_step}};
    \node[tag, below=1mm of s4] {+\texttt{tags}};
    \node[tag, below=1mm of s5] {train / eval\\ ($\approx$9000 / 1000)};
  \end{tikzpicture}
  \caption{数据流水线：从原始 Parquet 到训练/评测 JSONL 的五阶段流程}
  \label{fig:ch3-pipeline}
\end{figure}

\paragraph{（4）单条样本示例}
图~\ref{fig:ch3-sample} 摘录评测集首条样本（「将数组划分为 $k$ 个等和子集」）的 JSONL 记录。单条样本同时承载题干、参考解答、单测、函数签名、思维步骤与算法标签六类信息，共同支撑后续训练侧与推理侧的所有对照配置。少样本检索所需的 \texttt{strong\_tags} 不持久化，而在检索时与 3.3.2 节强标签常量取交集即时派生。

\begin{figure}[!htb]
\begin{footnotesize}
\begin{verbatim}
{
  "question":  "You are given an array of integers `arr` and an integer `k`.
                Partition the array into exactly `k` subsets ...",
  "solution":  "def subsets_with_sum(arr, k, target_sum, used): ...\n
                def can_partition(arr, k): ...",
  "test":      "from solution import can_partition\n
                def test_can_partition_with_valid_partition(): ...",
  "test_info": [{"function_name": "can_partition",
                 "function_declaration": "def can_partition(arr, k):",
                 "parameter_list": "(arr, k)", "docstring": "..."}, ...],
  "thought_step": "Step 1: [Understand] - ... Step 5: [Verify] - ...",
  "tags":      ["Backtracking", "Subset Sum", "Partition Problem"]
}
\end{verbatim}
\end{footnotesize}
\vspace{-2mm}
\caption{评测集首条样本的 JSONL 字段示意（内容经截断折行）}
\label{fig:ch3-sample}
\end{figure}
